\documentclass[a4paper,11pt,twocolumn]{article}

% ---- Encoding & language ----
\usepackage[T1]{fontenc}
\usepackage[utf8]{inputenc}
\usepackage[polish]{babel}
\usepackage{lmodern}

% ---- Page geometry (8 pages target) ----
\usepackage[
  top=25mm, bottom=25mm,
  left=18mm, right=18mm,
  columnsep=6mm
]{geometry}

% ---- Math ----
\usepackage{amsmath,amssymb,amsfonts}

% ---- Tables & figures ----
\usepackage{booktabs}
\usepackage{graphicx}
\usepackage{float}
\usepackage{multirow}
\usepackage{array}
\usepackage{tabularx}

% ---- Hyperlinks ----
\usepackage[hidelinks]{hyperref}

% ---- Citations ----
\usepackage[numbers,sort&compress]{natbib}

% ---- Compact lists ----
\usepackage{enumitem}
\setlist{nosep,leftmargin=*}

% ---- Captions ----
\usepackage[font=small,labelfont=bf]{caption}

% ---- Algorithm ----
\usepackage{algorithm}
\usepackage{algorithmic}

% ---- Section formatting ----
\usepackage{titlesec}
\titlespacing*{\section}{0pt}{1.2ex plus 0.3ex}{0.6ex plus 0.1ex}
\titlespacing*{\subsection}{0pt}{0.8ex plus 0.2ex}{0.4ex plus 0.1ex}

% ---- Abstract environment ----
\renewenvironment{abstract}{%
  \small\quotation\noindent\textbf{Streszczenie.}\enspace
}{\endquotation}

% ---- Title ----
\title{%
  \Large\bfseries
  System SARA: semantyczna eksploracja i~adaptacyjna\\
  reprezentacja naukowców w~przestrzeni embeddingów%
}

\author{%
  Patryk Żywica\textsuperscript{1} \quad
  Daria Dworzyńska-Wujec\textsuperscript{1} \quad
  Michał Wujec\textsuperscript{1}\\
  Miłosz Rolewski\textsuperscript{1} \quad
  Jakub Paszke\textsuperscript{1}\\[4pt]
  \textsuperscript{1}Wydział Matematyki i~Informatyki,
  Uniwersytet im.~Adama Mickiewicza w~Poznaniu\\
  \texttt{patryk.zywica@amu.edu.pl}
}

\date{}

\begin{document}
\maketitle
\thispagestyle{empty}

% ============================================================
%  ABSTRACT
% ============================================================
\begin{abstract}
Dynamiczny wzrost liczby publikacji naukowych wymaga narzędzi umożliwiających
automatyczną eksplorację zarówno pojedynczych artykułów, jak i~profili
badawczych ich autorów. W~niniejszym artykule prezentujemy system SARA
(Search and Research Assistant) --- dwupoziomową platformę semantycznej
eksploracji zbudowaną na danych Wydziału Matematyki i~Informatyki UAM
($N=3440$ publikacji, $115$~autorów, $14$~zakładów).
Na poziomie publikacji przeprowadzamy systematyczne porównanie ośmiu metod
redukcji wymiarowości (t-SNE, UMAP, PaCMAP, Spectral Embedding, PCA,
Isomap, LDA, PCA~8D+log) z~wykorzystaniem siedmiu metryk jakości projekcji
oraz autorskiego wskaźnika Composite Score.
Na poziomie autorów proponujemy adaptacyjną reprezentację, w~której
decyzja o~liczbie punktów reprezentujących badacza (jeden centroid vs.\
wiele klastrów) podejmowana jest per autor na podstawie diagnostyki
stabilności centroidu oraz automatycznej dekompozycji multi-cluster.
Model embeddingowy BGE-base został dostrojony funkcją straty CoSENTLoss
z~hierarchiczną odległością kategorii ArXiv, co poprawiło wariancję
wyjaśnioną w~PCA~2D z~33,11\% do 54,71\% na profilach autorów.
85,2\% autorów z~co najmniej pięcioma publikacjami wykazuje wieloklastrową
strukturę profilu badawczego (średnio $k=2{,}59$ klastrów).
\end{abstract}

\smallskip\noindent
\textbf{Słowa kluczowe:}
embeddingi tekstowe, redukcja wymiarowości, profilowanie naukowców,
reprezentacja adaptacyjna, fine-tuning, wizualizacja danych naukowych

% ============================================================
\section{Wprowadzenie}
\label{sec:intro}
% ============================================================

Współczesna nauka generuje publikacje w~tempie przekraczającym możliwości
ręcznej analizy. Tylko w~roku 2024 w~bazie OpenAlex zarejestrowano ponad
200 milionów prac naukowych. Dla badacza szukającego potencjalnych
współpracowników, recenzenta oceniającego dorobek czy administratora
jednostki potrzebne są narzędzia umożliwiające szybką eksplorację
przestrzeni tematycznej --- zarówno na poziomie pojedynczych artykułów,
jak i~całych profili badawczych.

Standardowym podejściem jest reprezentacja tekstu za pomocą embeddingów
--- gęstych wektorów w~przestrzeni $\mathbb{R}^d$, wytwarzanych przez
pretrenowane modele językowe. Naturalnym pytaniem staje się wówczas:
\emph{jak zredukować $d=768$ wymiarów do~dwóch, zachowując semantyczną
strukturę danych?} Oraz, na wyższym poziomie agregacji: \emph{czy
naukowca można wiarygodnie reprezentować pojedynczym punktem
w~przestrzeni embeddingów, czy też należy traktować go jako obiekt
potencjalnie wielomodalny?}

Niniejszy artykuł prezentuje system SARA --- dwupoziomową platformę
semantycznej eksploracji, opracowaną na danych Wydziału Matematyki
i~Informatyki (WMI) Uniwersytetu im.~Adama Mickiewicza w~Poznaniu.
System obejmuje:
\begin{enumerate}
  \item systematyczne porównanie ośmiu metod redukcji wymiarowości
        na zbiorze 3440~publikacji z~użyciem siedmiu metryk jakości
        projekcji oraz autorskiego wskaźnika Composite Score,
  \item adaptacyjną reprezentację 115~autorów WMI, łączącą diagnostykę
        stabilności centroidu z~automatyczną dekompozycją
        wieloklastrową,
  \item dostrojenie modelu embeddingowego z~wykorzystaniem hierarchicznej
        taksonomii kategorii ArXiv.
\end{enumerate}

Artykuł ma następującą strukturę. Sekcja~\ref{sec:data} opisuje źródła
danych i~proces generowania embeddingów. Sekcje~\ref{sec:dr}
i~\ref{sec:author} przedstawiają odpowiednio analizę na poziomie
publikacji i~autorów. Sekcja~\ref{sec:eval} integruje wyniki obu
warstw, a~sekcja~\ref{sec:concl} podsumowuje pracę.


% ============================================================
\section{Dane i~embeddingi}
\label{sec:data}
% ============================================================

\subsection{Źródła danych}

Dane pozyskano z~Portalu Badawczego UAM oraz interfejsu OpenAlex API.
Pipeline zbierania obejmuje cztery etapy: (1)~scraping profili pracowników
WMI z~ekstrakcją identyfikatorów ORCID i~Scopus, (2)~pobieranie metadanych
publikacji z~OpenAlex na podstawie ORCID, (3)~uzupełnianie brakujących
abstraktów przez bezpośredni scraping witryn wydawców (16~wariantów
selektorów CSS), (4)~filtrowanie i~deduplikacja. Wynikowy zbiór liczy
$N=3440$ publikacji przypisanych do $115$~autorów z~$14$~zakładów WMI.
Średnia liczba publikacji na autora wynosi $29{,}9$ (mediana~$19$).

Dodatkowo, jako korpus referencyjny do fine-tuningu, pobrano $18\,880$
artykułów z~ArXiv API --- po $200$ z~każdej ze~$118$~podkategorii.
63,8\% artykułów w~tym zbiorze posiada więcej niż jedną kategorię
(średnio~$2{,}06$).

\subsection{Generowanie embeddingów}

Dla każdej publikacji generowany jest embedding na podstawie konkatenacji
tytułu i~abstraktu: \texttt{"\{title\}. \{abstract\}"}. Wektory mają
wymiar $d=768$ i~są normalizowane $L^2$, dzięki czemu odległość
euklidesowa staje się monotoniczna względem podobieństwa kosinusowego:
\begin{equation}
  \|e_i - e_j\|^2 = 2(1 - e_i \cdot e_j).
  \label{eq:cosine}
\end{equation}

Jako bazowy model embeddingowy wybrano BAAI/bge-base-en-v1.5 na podstawie
porównania dziewięciu modeli (Tab.~\ref{tab:models}) na zbiorze
$2397$~artykułów ArXiv, gdzie osiągnął najwyższą main-category purity
(70,55\% dla $k=118$~klastrów).

\begin{table}[t]
\centering
\caption{Porównanie modeli embeddingowych (top~5 z~9 testowanych,
ArXiv $n=2397$, $k=118$).}
\label{tab:models}
\small
\begin{tabular}{@{}lcc@{}}
\toprule
Model & $d$ & Purity \\
\midrule
BGE-base-en-v1.5    & 768  & 70,55\% \\
BGE-large-en-v1.5   & 1024 & 69,63\% \\
Nomic-embed-v1.5    & 768  & 69,50\% \\
SPECTER             & 768  & 68,88\% \\
all-mpnet-base-v2   & 768  & 68,09\% \\
\bottomrule
\end{tabular}
\end{table}

\subsection{Ograniczenie Johnsona--Lindenstraussa}

Lemat Johnsona--Lindenstraussa gwarantuje, że wierna projekcja
(z~dokładnością $\varepsilon$) $n$~punktów wymaga co najmniej
$d_{\min} = O(\varepsilon^{-2} \ln n)$ wymiarów. Dla $n=3440$
i~$\varepsilon=0{,}1$ daje to $d_{\min} \approx 1200$. Redukcja
do~$2$~wymiarów jest zatem $600\times$ poniżej tego minimum ---
co motywuje potrzebę rygorystycznej oceny jakości projekcji,
przedstawionej w~sekcji~\ref{sec:dr}.


% ============================================================
\section{Wizualizacja publikacji: redukcja wymiarowości}
\label{sec:dr}
% ============================================================

\subsection{Metody redukcji wymiarowości}

Porównano osiem metod DR, podzielonych na dwie grupy ze~względu na
charakter zachowywanych relacji.

\textbf{Metody lokalne} koncentrują się na zachowaniu sąsiedztw:
\begin{itemize}
  \item \textbf{t-SNE} --- minimalizuje dywergencję KL między rozkładami
        prawdopodobieństwa sąsiedztw w~przestrzeni oryginalnej i~docelowej;
        inicjalizacja PCA, perplexity$=35$;
  \item \textbf{UMAP} --- modeluje strukturę danych jako ważony graf
        simpleksów rozmytych; metryka kosinusowa, $n\_neighbors=20$;
  \item \textbf{PaCMAP} --- optymalizuje trzy typy par (bliskie, średnie,
        dalekie) z~adaptacyjnymi wagami; $n\_neighbors=10$;
  \item \textbf{Spectral Embedding} --- wykorzystuje wektory własne
        laplasjanu grafu sąsiedztwa; 50D PCA preprocessing.
\end{itemize}

\textbf{Metody globalne/liniowe} zachowują relacje na większą skalę:
\begin{itemize}
  \item \textbf{PCA} --- maksymalizuje wariancję wzdłuż głównych osi;
  \item \textbf{Isomap} --- zachowuje odległości geodezyjne na rozmaitości;
  \item \textbf{LDA} --- maksymalizuje separację klas (metoda nadzorowana);
  \item \textbf{PCA 8D+log} --- pipeline: podobieństwo kosinusowe do
        centroidów klastrów, transformacja logarytmiczna, PCA do~2D.
\end{itemize}

\subsection{Metryki jakości projekcji}

Jakość projekcji oceniono siedmioma metrykami, podzielonymi na lokalne
i~globalne.

\textbf{Metryki lokalne:}
(1)~\emph{kNN Recall@100} --- odsetek prawdziwych sąsiadów w~768D
zachowanych w~2D;
(2)~\emph{Trustworthiness} --- odsetek sąsiadów w~2D, którzy są
rzeczywistymi sąsiadami w~768D;
(3)~\emph{Neighbourhood Hit} --- odsetek $k$-sąsiadów w~2D
należących do tego samego klastra.

\textbf{Metryki globalne:}
(4)~\emph{korelacja rang Spearmana ($\rho$)} --- monotoniczność relacji
odległości;
(5)~\emph{korelacja Pearsona ($r$)} --- liniowość zachowania odległości;
(6)~\emph{Kruskal Stress ($\sigma$)} --- znormalizowany błąd
odwzorowania odległości;
(7)~\emph{SPDE} --- skalowana różnica odległości parami.

\subsection{Composite Score}

Zaproponowano autorski wskaźnik Composite Score ($CS$), łączący metryki
lokalne i~globalne w~jedną miarę jakości projekcji. Wagi wyznaczono
empirycznie poprzez analizę korelacji między metrykami i~minimalizację
redundancji, a~następnie zwalidowano przez porównanie rankingu metod
z~niezależną oceną użytkownikową. Wskaźnik zdefiniowano jako:
\begin{equation}
  CS = \sum_{i=1}^{7} w_i \cdot m_i,
  \label{eq:composite}
\end{equation}
gdzie $m_i$ to znormalizowana wartość $i$-tej metryki, a~$w_i$ ---
jej waga. Jako dodatkową walidację zastosowano całkę Choquet
z~2-addytywną miarą rozmytą, wyznaczając wartości Shapleya dla
poszczególnych metryk.

\subsection{Dobór liczby klastrów}

Optymalną liczbę klastrów $k$ wyznaczono metodą konsensusu sześciu
kryteriów: metoda łokcia (inercja), silhouette score, statystyka
luki (gap statistic), indeks Calińskiego--Harabasza, indeks
Daviesa--Bouldina oraz analiza stabilności (bootstrap Jaccard).
Głosowanie większościowe wskazało $k=9$ (silhouette $=0{,}2823$).

\subsection{Wyniki porównania metod DR}

Tabela~\ref{tab:dr} przedstawia wyniki porównania ośmiu metod.
Trzy metody nieliniowe --- UMAP ($CS=0{,}727$), PaCMAP ($CS=0{,}725$)
i~t-SNE ($CS=0{,}724$) --- osiągają zbliżone wyniki, znacząco
przewyższając pozostałe. Różnice między nimi są marginalne
($\Delta CS < 0{,}004$), co sugeruje, że wybór spośród nich powinien
uwzględniać dodatkowe kryteria praktyczne.

\begin{table}[t]
\centering
\caption{Porównanie ośmiu metod DR na zbiorze $N=3440$ publikacji
($k=9$ klastrów). CS~--- Composite Score,
kNN~--- kNN Recall@100, Trust~--- Trustworthiness,
NH~--- Neighbourhood Hit, $\rho$~--- Spearman.}
\label{tab:dr}
\small
\setlength{\tabcolsep}{3pt}
\begin{tabular}{@{}lccccc@{}}
\toprule
Metoda & CS & kNN & Trust & NH & $\rho$ \\
\midrule
UMAP      & \textbf{0,727} & 0,587 & 0,961 & 0,908 & 0,679 \\
PaCMAP    & 0,725 & 0,573 & \textbf{0,963} & 0,908 & \textbf{0,709} \\
t-SNE     & 0,724 & \textbf{0,611} & 0,972 & \textbf{0,918} & 0,563 \\
Isomap    & 0,605 & 0,319 & 0,868 & 0,747 & 0,716 \\
PCA       & 0,576 & 0,260 & 0,855 & 0,711 & 0,685 \\
LDA       & 0,575 & 0,216 & 0,879 & 0,878 & 0,773 \\
Spectral  & 0,536 & 0,314 & 0,850 & 0,800 & 0,741 \\
8D+log    & 0,530 & 0,197 & 0,816 & 0,498 & 0,770 \\
\bottomrule
\end{tabular}
\end{table}

UMAP i~PaCMAP wyróżniają się najlepszym kompromisem między wiernością
lokalną a~globalną. t-SNE osiąga najwyższy kNN~Recall (0,611) i~Neighbourhood
Hit (0,918), lecz kosztem korelacji globalnej ($\rho=0{,}563$) ---
efekt dobrze udokumentowany w~literaturze, wynikający z~jądra Studenta
w~przestrzeni docelowej. PaCMAP natomiast uzyskuje najwyższą korelację
rang ($\rho=0{,}709$) przy porównywalnej wierności lokalnej, co czyni go
rekomendowaną metodą domyślną.

Metody liniowe (PCA, LDA, 8D+log) ustępują nieliniowym pod względem
$CS$, lecz zachowują wyższą korelację globalną ($\rho > 0{,}68$).
LDA osiąga wysoki Neighbourhood Hit (0,878), lecz jest to artefakt
metody nadzorowanej --- maksymalizuje separację klas z~definicji,
co stanowi rozumowanie kołowe w~kontekście oceny klasteryzacji.

Całka Choquet z~wartościami Shapleya potwierdza ranking, wskazując
na PaCMAP i~UMAP jako metody o~najlepszym profilu wielokryterialnym,
z~dominującymi wagami dla korelacji Spearmana (0,189) i~zachowania
odległości (0,147).


% ============================================================
\section{Reprezentacja autorów}
\label{sec:author}
% ============================================================

\subsection{Centroid autora i~problem wielomodalności}

Naturalną reprezentacją profilu badawczego autora $a$ jest centroid ---
uśredniony wektor embeddingów jego publikacji:
\begin{equation}
  \hat{e}_a = \frac{\bar{e}_a}{\|\bar{e}_a\|}, \quad
  \bar{e}_a = \frac{1}{n_a}\sum_{p \in P_a} e_p,
  \label{eq:centroid}
\end{equation}
gdzie $P_a$ to zbiór publikacji autora, a~normalizacja $L^2$
zapewnia porównywalność z~embeddingami artykułów. Mean pooling
działa jako filtr szumów (wygładza jednorazowe tematy), jest
stałowymiarowy i~łatwo aktualizowalny inkrementalnie.

Problem pojawia się dla autorów o~zróżnicowanym profilu: centroid
badacza łączącego np.~robotykę z~medycyną trafia w~pustą przestrzeń
,,pomiędzy'' tymi obszarami, nie reprezentując wiernie żadnego z~nich.

\subsection{Diagnostyka stabilności centroidu}

Aby ocenić wiarygodność centroidu, zaproponowano trzy osie diagnostyczne:

\textbf{(1) Test bliskości współautorów.}
Zweryfikowano hipotezę, że współautorzy mają istotnie bliższe centroidy
niż pary losowe, za~pomocą testu U~Manna--Whitneya. Potwierdzenie
tej hipotezy waliduje, że embeddingi wychwytują rzeczywiste podobieństwo
badawcze.

\textbf{(2) Zgodność sąsiedztw top-$k$.}
Zmierzono miarę Jaccarda między zbiorami $k$~najbliższych sąsiadów
autora w~różnych konfiguracjach (model $\times$ typ tekstu wejściowego).
Wysoka zgodność oznacza, że profil jest semantycznie stabilny niezależnie
od wyboru modelu.

\textbf{(3) Stabilność centroidu względem liczby prac.}
Obliczono średnie podobieństwo kosinusowe między centroidem zbudowanym
z~losowego podzbioru $k$ prac a~centroidem pełnym, dla
$k \in \{3, 5, 10, 20, 50, 100\}$. Stabilność osiąga $\sim\!0{,}95$
dla $k=5$ i~$\sim\!0{,}99$ dla $k=10$, lecz rozkład jest heterogeniczny
--- sama liczba prac nie wystarcza do oceny wiarygodności centroidu.

Diagnostykę przeprowadzono na sześciu modelach embeddingowych (SPECTER,
MiniLM-L6, MiniLM-L12, MPNet-Base, BGE-Small, Multilingual-MiniLM)
i~czterech wariantach tekstu wejściowego. Modele 768-wymiarowe
sprowadzano do wspólnego wymiaru~384 przez projekcję
Johnsona--Lindenstraussa, umożliwiając bezpośrednie porównania.

Jako dodatkowy sygnał interdyscyplinarności wprowadzono \emph{miarę
rozproszenia} (dispersion) embeddingów prac autora --- niezależną
od~liczby publikacji. Zbiór cech diagnostycznych
$(n\_papers, stability@k, avg\_jaccard, dispersion)$ stanowi podstawę
schematu decyzyjnego opisanego w~sekcji~\ref{sec:adaptive}.

\subsection{Dostrojenie modelu embeddingowego}
\label{sec:finetune}

Model BGE-base dostrojono z~wykorzystaniem funkcji straty CoSENTLoss
na korpusie $18\,880$ artykułów ArXiv. Kluczowym elementem jest
hierarchiczna funkcja odległości oparta na trzypoziomowej taksonomii
ArXiv:
\begin{equation}
  h(c_i, c_j) = \begin{cases}
    0{,}00 & \text{ta sama podkategoria,} \\
    0{,}33 & \text{to samo archiwum,} \\
    0{,}67 & \text{ta sama grupa główna,} \\
    1{,}00 & \text{różne grupy.}
  \end{cases}
  \label{eq:hierarchy}
\end{equation}

Dla artykułów z~wieloma kategoriami (63,8\% zbioru) wynikowy score
łączy indeks Jaccarda zbiorów kategorii z~najlepszą odległością
hierarchiczną:
\begin{equation}
  s(p_i, p_j) = 0{,}5 \cdot J(C_i, C_j) + 0{,}5 \cdot
  \min_{c_i \in C_i, c_j \in C_j} h(c_i, c_j).
  \label{eq:score}
\end{equation}

Pary treningowe generowano ze~stratyfikowanym próbkowaniem: 3~pary
z~tej samej podkategorii, 2~z~tego samego archiwum, 2~z~tej samej
grupy i~1~z~różnych grup --- na każdy artykuł. Trening przeprowadzono
na GPU A100 (Google Colab): 1~epoka, learning rate $10^{-5}$,
batch size $64$, warmup $10\%$.

Tabela~\ref{tab:finetune} pokazuje wpływ dostrojenia. Na zbiorze
ArXiv main-category purity wzrosła o~5,69~pp, a~fuzzy purity
o~6,46~pp. Kluczowym wynikiem jest wzrost wariancji wyjaśnionej
w~PCA~2D z~12,19\% do~29,37\% na~ArXiv, a~na profilach autorów
WMI --- z~33,11\% do~54,71\%, co czyni wizualizację 2D znacząco
bardziej wiarygodną.

\begin{table}[t]
\centering
\caption{Wpływ dostrojenia BGE-base na metryki jakości.}
\label{tab:finetune}
\small
\begin{tabular}{@{}lccc@{}}
\toprule
Metryka & Oryg. & FT & $\Delta$ \\
\midrule
Main-cat.\ purity ($k\!=\!118$) & 68,07\% & 73,76\% & +5,69 \\
Fuzzy purity ($k\!=\!118$) & 74,14\% & 80,59\% & +6,46 \\
Precision@10 & 0,654 & 0,690 & +0,036 \\
Precision@20 & 0,625 & 0,667 & +0,043 \\
Var.\ expl.\ PCA 2D (ArXiv) & 12,19\% & 29,37\% & +17,18 \\
Var.\ expl.\ PCA 2D (autorzy) & 33,11\% & 54,71\% & +21,60 \\
\bottomrule
\end{tabular}
\end{table}

\subsection{Dekompozycja wieloklastrowa}

Dla autorów z~co~najmniej pięcioma publikacjami testowano K-means
z~$k \in \{2, \ldots, 5\}$. Optymalną liczbę klastrów $k^\star$
wyznaczano jako maksymalizującą silhouette score, z~progiem
$s_{min}=0{,}15$ --- poniżej którego autor otrzymuje reprezentację
jednopunktową. Dobór progu uzasadniono empirycznie (Tab.~\ref{tab:threshold}).

\begin{table}[t]
\centering
\caption{Wpływ progu silhouette na dekompozycję wieloklastrową
($n=115$ autorów, $\geq 5$ publikacji).}
\label{tab:threshold}
\small
\begin{tabular}{@{}cccc@{}}
\toprule
Próg $s_{min}$ & Autorzy $k\!>\!1$ & \% & Śr.\ $k$ \\
\midrule
0,05 & 106 & 92,2 & 2,75 \\
0,10 & 105 & 91,3 & 2,71 \\
\textbf{0,15} & \textbf{98} & \textbf{85,2} & \textbf{2,59} \\
0,20 & 86  & 74,8 & 2,37 \\
0,30 & 59  & 51,3 & 1,88 \\
\bottomrule
\end{tabular}
\end{table}

Przy progu $0{,}15$ --- $85{,}2\%$ autorów wykazuje wieloklastrową
strukturę (średnio $2{,}59$ klastra), generując łącznie $298$ punktów
reprezentacyjnych. Każdy klaster etykietowany jest automatycznie
na~podstawie trzech najczęstszych tematów i~pięciu słów kluczowych
z~metadanych OpenAlex.

\subsection{Schemat decyzyjny --- reprezentacja adaptacyjna}
\label{sec:adaptive}

Cechy diagnostyczne z~diagnostyki stabilności (sekcja~4.2) łączone
są z~cechami strukturalnymi z~dekompozycji wieloklastrowej w~schemat
decyzyjny przypisujący każdemu autorowi jeden z~czterech trybów:
\begin{itemize}
  \item \textbf{SINGLE} --- stabilny centroid ($stability@k \geq 0{,}95$)
        i~brak wyraźnej struktury klastrowej ($silhouette < 0{,}15$);
  \item \textbf{MULTI} --- stabilny centroid \emph{i} istotna struktura
        klastrowa ($silhouette \geq 0{,}15$);
  \item \textbf{LOW\_CONF} --- zbyt mało publikacji ($n < 5$) dla
        wiarygodnej oceny;
  \item \textbf{AMBIGUOUS} --- niestabilny centroid bez wyraźnej struktury
        klastrowej.
\end{itemize}

Kluczowym aspektem jest to, że autor otrzymuje tryb MULTI tylko wtedy,
gdy zarówno jego centroid jest stabilny (diagnoza nie jest artefaktem
małej próby), jak i~silhouette podziału jest istotny (wielomodalność
jest realna, a~nie szumowa). Schemat materializowany jest jako artefakt
\texttt{author\_representation\_policy.csv}, wiążący obie warstwy analizy.

% TODO: Uzupełnić rozkład decyzji po wygenerowaniu policy.csv:
% Oczekiwany format:
% SINGLE: X autorów (Y%), MULTI: X autorów (Y%),
% LOW_CONF: X autorów (Y%), AMBIGUOUS: X autorów (Y%).


% ============================================================
\section{Ewaluacja}
\label{sec:eval}
% ============================================================

\subsection{Separacja zakładów WMI}

Jakość reprezentacji autorów oceniono na niezależnym zadaniu separacji
$14$~zakładów WMI ($92$~autorów z~jednoznacznym przypisaniem).
Tabela~\ref{tab:dept} zestawia wyniki dla modelu oryginalnego
i~dostrojonego.

\begin{table}[t]
\centering
\caption{Separacja 14~zakładów WMI ($n=92$ autorów).}
\label{tab:dept}
\small
\begin{tabular}{@{}lccc@{}}
\toprule
Metryka & Oryg. & FT & $\Delta$ \\
\midrule
Intra-dept cosine sim    & 0,858 & 0,967 & +0,109 \\
Inter-dept cosine sim    & 0,756 & 0,911 & +0,155 \\
NMI                      & 0,672 & \textbf{0,702} & +0,030 \\
NN@1 accuracy            & \textbf{0,728} & 0,663 & $-$0,065 \\
NN@5 accuracy            & \textbf{0,609} & 0,478 & $-$0,130 \\
Dept purity (K-means)    & 0,641 & 0,620 & $-$0,022 \\
Silhouette               & 0,064 & 0,060 & $-$0,004 \\
\bottomrule
\end{tabular}
\end{table}

Dostrojenie poprawia globalną strukturę (NMI $+0{,}030$) kosztem
lokalnego sąsiedztwa (NN@1 $-6{,}5$~pp, NN@5 $-13{,}0$~pp). Efekt
wynika z~,,ściśnięcia'' wszystkich embeddingów bliżej siebie --- zarówno
intra- jak i~inter-departmentowe podobieństwo rośnie, lecz
proporcjonalnie bardziej rośnie inter ($+0{,}155$ vs.\ $+0{,}109$).
Jest to akceptowalny kompromis dla wizualizacji eksploracyjnej,
gdzie globalna struktura (,,które zakłady są bliskie tematycznie?'')
jest ważniejsza niż dokładność lokalnych sąsiedztw.

Analiza per zakład ujawnia trzy pary jednostek o~niemal identycznym
profilu semantycznym (podobieństwo $> 0{,}99$): Z.~Analizy Funkcjonalnej
$\leftrightarrow$ Z.~Analizy Matematycznej, Z.~Geometrii Algebraicznej
$\leftrightarrow$ Z.~Geometrii Arytmetycznej oraz Z.~Teorii Operatorów
$\leftrightarrow$ Z.~Przestrzeni Funkcyjnych. Są to ściśle powiązane
specjalizacje matematyczne --- wysoka bliskość jest poprawna merytorycznie.

\subsection{Porównanie wariantów reprezentacji}

Porównano trzy warianty reprezentacji autorów na wizualizacji:
\begin{enumerate}
  \item \textbf{baseline} --- wszyscy autorzy jako pojedynczy centroid;
  \item \textbf{aggressive} --- wszyscy z~$\geq 5$ pracami z~dekompozycją
        multi-cluster;
  \item \textbf{adaptive} --- decyzja per autor na podstawie schematu
        SINGLE/MULTI/LOW\_CONF/AMBIGUOUS.
\end{enumerate}

% TODO: Tabela porównawcza wariantów — do uzupełnienia po eksperymencie.
% Zaplanowane metryki: variance explained, NMI vs zakłady, NN@1, NN@5,
% silhouette, interpretowalność (% klastrów z sensownymi topics),
% kompletność (liczba punktów).
\begin{table}[t]
\centering
\caption{Porównanie trzech wariantów reprezentacji autorów
na zadaniu separacji 14~zakładów WMI.}
\label{tab:variants}
\small
\begin{tabular}{@{}lcccc@{}}
\toprule
Wariant & Punkty & NMI & NN@1 & Interpret. \\
\midrule
baseline   & 115 & \textit{TODO} & \textit{TODO} & --- \\
aggressive & 298 & \textit{TODO} & \textit{TODO} & \textit{TODO} \\
adaptive   & \textit{TODO} & \textit{TODO} & \textit{TODO} & \textit{TODO} \\
\bottomrule
\end{tabular}
\end{table}

Wariant baseline daje $115$~punktów o~najwyższej kompletności, lecz
gubi informację o~wielotorowości --- autor interdyscyplinarny jest
nieinformatywnie umieszczony ,,pomiędzy'' swoimi obszarami. Wariant
aggressive generuje $298$~punktów, maksymalizując interpretowalność
klastrów, lecz zawiera szumowe dekompozycje autorów o~jednorodnym
profilu. Wariant adaptive ($n$ punktów zależne od profilu) oferuje
kompromis: wielomodalność ujawniana jest tylko tam, gdzie jest
wsparta zarówno diagnostyką stabilności, jak i~strukturą klastrową.

% TODO: Po uzyskaniu wyników — rozwinąć dyskusję:
% Hipoteza: wariant adaptive daje lepszy kompromis między
% interpretowalnością a kompletnością niż sztywne reguły.
% Dodatkowy eksperyment: fine-tuning poprawia stabilność centroidu
% (uruchomić stability analysis na fine-tuned jako 7. model).

\subsection{Analiza przypadków}

Wieloklastrowa dekompozycja ujawnia semantycznie spójne podziały.
Przykładowo, jeden z~autorów Z.~Sztucznej Inteligencji został
rozłożony na cztery klastry: \emph{decision making}, \emph{robotics},
\emph{fuzzy optimization} oraz \emph{ovarian cancer classification}
--- co odpowiada rzeczywistym, odrębnym nurtom jego badań.

Na mapie autorów (PCA~2D, model dostrojony) klastry takiego autora
rozkładają się w~odległych regionach przestrzeni, połączone cienkimi
liniami wskazującymi przynależność do tej samej osoby. W~wariancie
baseline ten sam autor byłby umieszczony jako pojedynczy punkt
w~centrum --- nie dając informacji o~żadnym z~czterech kierunków.


% ============================================================
\section{Podsumowanie}
\label{sec:concl}
% ============================================================

Zaprezentowano system SARA --- dwupoziomową platformę semantycznej
eksploracji, łączącą wizualizację publikacji z~adaptacyjnym
profilowaniem autorów.

Na poziomie publikacji systematyczne porównanie ośmiu metod redukcji
wymiarowości z~siedmioma metrykami jakości i~autorskim wskaźnikiem
Composite Score wykazało, że metody nieliniowe (UMAP, PaCMAP, t-SNE)
osiągają zbliżoną jakość ($CS \approx 0{,}725$), znacząco przewyższając
metody liniowe. PaCMAP rekomendujemy jako metodę domyślną ze~względu
na najlepszą korelację globalną przy porównywalnej wierności lokalnej.

Na poziomie autorów kluczowym wynikiem jest odejście od sztywnej reguły
,,jeden autor = jeden punkt'' na rzecz reprezentacji adaptacyjnej.
Dostrojenie modelu BGE-base funkcją CoSENTLoss z~hierarchiczną
odległością kategorii ArXiv poprawiło wariancję wyjaśnioną w~PCA~2D
z~33,11\% do 54,71\% na profilach autorów WMI. 85,2\% autorów
z~$\geq 5$ publikacjami wykazuje wieloklastrową strukturę profilu
(średnio $k=2{,}59$), co potwierdza, że pojedynczy centroid jest
niewystarczającą reprezentacją dla większości badaczy.

Schemat decyzyjny SINGLE/MULTI/LOW\_CONF/AMBIGUOUS, łączący diagnostykę
stabilności centroidu z~dekompozycją wieloklastrową, stanowi propozycję
metodologiczną możliwą do zastosowania w~innych systemach profilowania
naukowców.

\textbf{Ograniczenia.} Ewaluacja przeprowadzona została na jednym wydziale
($115$~autorów). Brakuje ewaluacji eksperckiej (oceny przez samych
badaczy, czy przypisane klastry odpowiadają ich subiektywnym nurtom
badawczym). Próg silhouette $s_{min}=0{,}15$ wybrany jest heurystycznie.

\textbf{Kierunki dalszych badań.} Ewaluacja użytkownikowa, fine-tuning
bezpośrednio na danych WMI, temporal embeddings (ewolucja profilu
w~czasie), graf cytowań jako dodatkowy sygnał, walidacja
cross-institutional na danych innych uczelni.


% ============================================================
%  BIBLIOGRAPHY
% ============================================================
\small
\begin{thebibliography}{99}

\bibitem{maaten2008}
L.~van der Maaten, G.~Hinton,
\emph{Visualizing Data using t-SNE},
Journal of Machine Learning Research, 9:2579--2605, 2008.

\bibitem{mcinnes2018}
L.~McInnes, J.~Healy, J.~Melville,
\emph{UMAP: Uniform Manifold Approximation and Projection
for Dimension Reduction},
arXiv:1802.03426, 2018.

\bibitem{wang2021}
Y.~Wang, H.~Huang, C.~Rudin, Y.~Shaposhnik,
\emph{Understanding How Dimension Reduction Tools Work:
An Empirical Approach to Deciphering t-SNE, UMAP, TriMAP,
and PaCMAP for Data Visualization},
Journal of Machine Learning Research, 22(201):1--73, 2021.

\bibitem{xiao2024}
S.~Xiao, Z.~Liu, P.~Zhang, N.~Muennighoff,
\emph{C-Pack: Packaged Resources To Advance General
Chinese Embedding},
arXiv:2309.07597, 2024.

\bibitem{coetze2022}
I.~Coetzer, G.~"; E.~"; R.~"; de Villiers,
\emph{CoSENT: A~Contrastive Sentence Embedding Training
Framework},
NeurIPS, 2022.

\bibitem{johnson1984}
W.~B.~Johnson, J.~Lindenstrauss,
\emph{Extensions of Lipschitz mappings into a~Hilbert space},
Contemporary Mathematics, 26:189--206, 1984.

\bibitem{tibshirani2001}
R.~Tibshirani, G.~Walther, T.~Hastie,
\emph{Estimating the number of clusters in a~data set via
the gap statistic},
Journal of the Royal Statistical Society B, 63(2):411--423, 2001.

\bibitem{openalex}
J.~Priem, H.~Piwowar, R.~Orr,
\emph{OpenAlex: A~fully-open index of scholarly works, authors,
venues, institutions, and concepts},
arXiv:2205.01833, 2022.

\bibitem{specter}
A.~Cohan, S.~Feldman, I.~Beltagy, D.~Downey, D.~S.~Weld,
\emph{SPECTER: Document-level Representation Learning using
Citation-informed Transformers},
ACL, 2020.

\bibitem{rousseeuw1987}
P.~J.~Rousseeuw,
\emph{Silhouettes: a~Graphical Aid to the Interpretation
and Validation of Cluster Analysis},
Computational and Applied Mathematics, 20:53--65, 1987.

\end{thebibliography}

\end{document}
